\chapter{Preliminaries}

Unless stated otherwise, we consider unoriented simple connected graphs $G=(V,E)$. Then we can replace each unoriented edge $uv$ with two oriented edges $\vv{uv}$, $\vv{vu}$, called \emph{darts}. By a \emph{dart-set} $D(S)$ of $S\subseteq V$ we mean the set of all darts with initial vertex in the set $S$. Machimal degree of a graph is denoted by $\Delta(G)$. A graph is called \emph{$r$-regular}, if each its vertex is of degree $r$. \emph{Cubic} graph is a $3$-regular graph. A \emph{dangling edge} is an edge with only one endpoint in a vertex of the graph. A \emph{$k$-pole} is a graph containing $k$ dangling edges.

A \emph{girth} is a length of the shortest cycle in a graph. A set of cycles covering all vertices exactly once is called a \emph{$2$-factor}. The smallest possible number of cycles on odd vertices in a $2$-factor of a graph is called \emph{oddness}.

A \emph{bridge} is an edge, whose removal disconnects the graph. Similarly, an \emph{edge-cut} is a set of edges disconnecting graph after its removal. An \emph{$uv$-edge-separator} is an edge-cut, for which $u$ and $v$ are in different components. An \emph{edge-connectivity} is a minimal number of edges in the edge-cut of the graph. The Menger's theorem \cite{menger} relates the size of the smallest $uv$-edge-separator to the largest number of disjoint $uv$-paths. We need its oriented version, so we state it below as Theorem \ref{th:oriented_menger}. A \emph{$2$-edge-connected graph} can be also called \emph{bridgeless}. For a \emph{cyclic edge-connectivity}, we also require the components to contain a cycle -- so they cannot be trees.

\begin{theorem}[Menger's theorem for oriented graphs] \emph{\cite[p. 149]{oriented_menger}}
	Let $u, v$ be the vertices of an \emph{oriented} multigraph $G$. Then there echist $k$ edge-disjoint oriented $uv$-paths if and only if for any set $S\subseteq V(G)\setminus\{v\}$ containing $u$, there are at least $k$ oriented edges from $S$ to $V(G)\setminus S$. \label{th:oriented_menger}
\end{theorem}

An \emph{edge-colouring} of the graph is an assignment from the set of edges $E$ to the set of colours $C$. The edge-colouring is \emph{proper}, if no adjacent edges share a colour. The minimal $k$, such that a graph has a proper edge-colouring using $k$ colours, is a \emph{chromatic index $\chi'(G)$}. The chromatic index is bounded by $\Delta(G)\leq\chi'(G)\leq\Delta(G)+1$ \cite{vizing}, and so can obtain only two values. Therefore, graphs with $\chi'(G)=\Delta(G)$ are called ``class one'', or ``colourable'' graphs, and the others are called ``class two'', or ``uncolourable'' graphs. A snark is an uncolourable cubic graph, which is cyclically $4$-edge-connected and has girth at least $5$.

A \emph{Petersen graph} (in Figure \ref{fig:petersen}) is the smallest snark. It is \emph{edge-transitive}, which means, that for any pair of edges $e_1, e_2$, there echists a bijection $\sigma\colon V\to V$ that preserves adjacency (i.e. $uv$ is an edge if and only if $\sigma(u)\sigma(v)$ is an edge) and $\sigma(e_1)=e_2$.
\begin{figure}[h!]
	\centering
	\input images/petersen.tex
	\caption{Petersen graph}
	\label{fig:petersen}
\end{figure}

\emph{Blanuša snarks} (in Figure \ref{fig:blanusa}) are the smallest snarks except the Petersen graph. They are a part of two infinite classes of snarks, called \emph{generalised Blanuša snarks}, constructed from three types of $4$-poles depicted in Figure \ref{fig:networks}. The \emph{$n$-th generalised Blanuša snark of type I $B_n^{(1)}$} is constructed from $n$ copies of the $4$-pole $B$ and one copy of the $4$-pole $I$ by identifying consecutive $a$ with $\overline{a}$ and $b$ with $\overline{b}$. Similarly, we obtain the \emph{$n$-th generalised Blanuša snark of type II $B_n^{(2)}$} from $n-1$ copies of the $4$-pole $B$ and one copy of the $4$-pole $X$. For both classes, the first graph in the class is the Petersen graph and the second one is the corresponding Blanuša snark.
\begin{figure}[h!]
	\begin{subfigure}{0.4\textwidth}
	\centering
	\input images/blanusa_I.tex
	\caption{first}
	\end{subfigure}
	\begin{subfigure}{0.4\textwidth}
	\centering
	\input images/blanusa_II.tex
	\caption{second}
	\end{subfigure}
	\caption{Blanuša snarks}
	\label{fig:blanusa}
\end{figure}

\begin{figure}[h!]
	\centering
	\begin{subfigure}{0.3\textwidth}
	\centering
	\input images/blanusa_network_B.tex
	\caption{$B$}
	\end{subfigure}
	\begin{subfigure}{0.2\textwidth}
	\centering
	\input images/blanusa_network_I.tex
	\caption{$I$}
	\end{subfigure}
	\begin{subfigure}{0.3\textwidth}
	\centering
	\input images/blanusa_network_X.tex
	\caption{$X$}
	\end{subfigure}
	\caption{Three $4$-poles used to construct generalised Blanuša snarks}
	\label{fig:networks}
\end{figure}
